\textbf{Hypotheses and tests (registered in \texttt{proposal.md} before the runs).} \textbf{H1:} at 128 points the FNO has lower test error than every baseline; test: Welch $t$-test on the five per-seed errors, $\alpha=0.05$, plus a gain of at least 10\% over the best baseline. \textbf{H2:} zero-shot on the 256-point grid, the FNO error ratio (256-point over 128-point error, mean over seeds) is below 1.5 while the CNN ratio exceeds 1.5. \textbf{H3:} FNO error saturates in the number of modes: 8 to 16 modes changes the mean error by under 10\%, and 4 modes is clearly worse (three seeds). \textbf{H4:} the FNO beats DeepONet at 100, 200 and 400 training pairs (three seeds).

\textbf{Data.} Initial conditions are random Fourier series, $a(x)=\sum_{k=1}^{63}s_k(a_k\cos2\pi kx+b_k\sin2\pi kx)$, $a_k,b_k\sim\mathcal N(0,1)$, $s_k=25\sqrt2/((2\pi k)^2+25)$, a Gaussian-measure family in the spirit of the FNO Burgers benchmark. The solution at $t=1$ is computed with an ETDRK4 pseudo-spectral solver~\cite{kassam2005fourth} on 512 points with $\Delta t=10^{-3}$ and 2/3 dealiasing, then subsampled to 128, 256 or 512 points. Solver self-checks (\texttt{method/data.py}) halve $\Delta t$ and re-solve on the 256- and 128-point grids; both change the solution by amounts far below every model error reported here, but they were run outside the run registry, so we quote no value from them. The logged counterpart is the interpolation floor recorded with every run: Fourier-interpolating the 128-point test solutions to 256 points reproduces the 256-point solutions to a relative error of \rhval{main/fno/burgers_nu0.01/interp_floor_256/mean}. All grids therefore resolve the solution and the solver is not a confound. Per seed we draw 600 initial conditions (generator seed $1000+s$) and split them 400/100/100 into train/validation/test. The seed changes the data, the initialisation and the batch order; the main comparison uses seeds 0--4, sweeps use seeds 0--2.

\textbf{Training and tuning.} All systems use the loss, optimiser, schedule, batch size and 100-epoch budget of Section~\ref{sec:method}. The only tuned hyperparameter is the learning rate, chosen per system from $\{3\!\times\!10^{-4},10^{-3},3\!\times\!10^{-3},10^{-2}\}$ on a separate tuning seed (100, disjoint data) by validation error, with no test data; Table~\ref{tab:tune} lists all sixteen points. The grid originally had three values. It was extended to $10^{-2}$ after the FNO's best rate sat at the upper edge, and DeepONet and the MLP were run at $10^{-2}$ at that time. The CNN's best rate sat at the same edge, but its $10^{-2}$ point was run only in the revision of this paper, after the main runs; it diverges, as DeepONet does at that rate (validation errors of \rhval{tune/tune-cnn@lr=0.01/burgers_nu0.01/rel_l2_val/mean:2} and \rhval{tune/tune-deeponet@lr=0.01/burgers_nu0.01/rel_l2_val/mean:2}, i.e.\ no better than predicting zero), so the selection is unchanged. The selected rates are $10^{-2}$ (FNO), $3\times10^{-3}$ (DeepONet and CNN) and $10^{-3}$ (MLP, whose validation errors at $10^{-3}$ and $3\times10^{-3}$, \rhval{tune/tune-mlp@lr=0.001/burgers_nu0.01/rel_l2_val/mean} and \rhval{tune/tune-mlp@lr=0.003/burgers_nu0.01/rel_l2_val/mean}, are nearly equal). The FNO's best rate is still at the upper edge of the extended grid. Two early DeepONet tuning runs crashed while writing their output file, before any metric was recorded; they were rerun and are listed as failed in the registry. Model sizes are in Table~\ref{tab:main}; the FNO uses 4 layers, width 32 and 16 modes.

\begin{table}[t]\centering\caption{Learning-rate tuning: validation relative $L_2$ at 128 points on the tuning seed (one run per cell, 100 epochs); bold marks the selected rate.}\label{tab:tune}
\begin{tabular}{lcccc}
\toprule
System & $3\!\times\!10^{-4}$ & $10^{-3}$ & $3\!\times\!10^{-3}$ & $10^{-2}$ \\
\midrule
FNO & 0.0132 & 0.0078 & 0.0062 & \textbf{0.0053} \\
DeepONet & 0.4147 & 0.3175 & \textbf{0.2056} & 1.0027 \\
MLP & 0.1027 & \textbf{0.0940} & 0.0941 & 0.2134 \\
CNN & 0.2512 & 0.2263 & \textbf{0.2175} & 1.0000 \\
\bottomrule
\end{tabular}
\end{table}

\textbf{Metrics, hardware.} Mean per-sample relative $L_2$ on the 100 test pairs at 128 (training grid), 256 and 512 points, with the zero-shot rules of Section~\ref{sec:method}. Everything ran on a shared laptop CPU with 2 torch threads, one or two processes at a time; per-run wall-clock training time is reported in the main table, and each run stayed under six minutes. All tables report mean $\pm$ standard deviation over seeds.
